% Generated draft: requires mathematical and source review.
\documentclass[11pt]{article}
\usepackage[margin=25mm]{geometry}
\usepackage{fontspec}
\setmainfont{Latin Modern Roman}
\usepackage{amsmath,amssymb,mathtools,unicode-math}
\setmathfont{Latin Modern Math}
\usepackage{enumitem,needspace,xurl,hyperref}
\setlength{\parindent}{0pt}
\setlength{\parskip}{5pt}
\setlength{\emergencystretch}{3em}
\newcommand{\N}{\mathbb{N}}
\newcommand{\Z}{\mathbb{Z}}
\newcommand{\Q}{\mathbb{Q}}
\newcommand{\R}{\mathbb{R}}
\newcommand{\C}{\mathbb{C}}
\newcommand{\F}{\mathbb{F}}
\newcommand{\A}{\mathbb{A}}
\newcommand{\PP}{\mathbb{P}}
\newcommand{\E}{\mathbb{E}}
\newcommand{\eps}{\varepsilon}
\newcommand{\dd}{\,\mathrm{d}}
\begin{document}
\section*{UnsolvedMath ID: 701073}
\textbf{Source problem code:} catalogue entry
\subsection{Definitions and mathematical statement}
Let \(p\) range over all primes, and let \(P\) be a finite group of order \(p^m\), where \(m\in\Z_{\geq 0}\). Write uniquely
\[
m=2n+e,\qquad n\in\Z_{\geq 0},\quad e\in\{0,1\}.
\]
Let \(k(P)\) denote the number of conjugacy classes of \(P\). The supplied source attributes to P. Hall the assertion that there is an integer \(a(P)\geq 0\) such that
\[
k(P)=n(p^2-1)+p^e+a(P)(p^2-1)(p-1).
\]
This integer \(a(P)\) is called the \emph{abundance} of \(P\).

Define the lower central series by \(\gamma_1(P)=P\) and \(\gamma_{i+1}(P)=[\gamma_i(P),P]\). Let \(c(P)\) be the least nonnegative integer \(c\) such that \(\gamma_{c+1}(P)=1\), and define the \emph{coclass} by
\[
r(P)=m-c(P),\qquad |P|=p^{c(P)+r(P)}.
\]
This convention assigns class \(0\) to the trivial group; the source does not specify its convention for that boundary case. For \(s\in P\), define
\[
C_P(s)=\{x\in P:xs=sx\}.
\]

\begin{enumerate}
\item[(a)] Does there exist a function \(g:\Z_{\geq 0}\to\Z_{\geq 0}\) such that, for every prime \(p\) and every finite \(p\)-group \(P\),
\[
r(P)\leq g(a(P))?
\]
Thus the bound must be independent of both \(p\) and \(|P|\).
\item[(b)] Does there exist a function \(f:\Z_{\geq 0}\to\Z_{\geq 0}\) such that, for every prime \(p\) and every finite \(p\)-group \(P\), there exists \(s\in P\) satisfying
\[
|C_P(s)|\leq p^{f(a(P))}?
\]
Here \(f\) must be independent of \(p\) and \(|P|\), although \(s\) may depend on \(P\).
\end{enumerate}

The source states that abundance \(0\) implies coclass \(1\), and reports that all examples known to its authors with abundance \(1\) have coclass at most \(3\). The first assertion requires a boundary qualification: a group of order \(p\) has abundance \(0\) and coclass \(0\). No restriction excluding this case is supplied. The statement about known abundance-\(1\) examples is a source-specific observation, not a universal assertion.

For part (b), the source states that \(f(0)=2\) is available and tentatively suggests \(f(1)=3\). The latter value is not presented as established.

The source also attributes to A. Jaikin-Zapirain a result asserting the existence of a function \(h\), depending only on the prime \(p\) and the abundance \(a\), such that
\[
|P|\leq p^{h(p,a(P))}
\]
for every finite \(p\)-group \(P\). The notation \(h\) distinguishes this supplied background bound from the function sought in part (b).
\subsection{Short English statement}
For a finite \(p\)-group, define its abundance by the supplied formula for its number of conjugacy classes. Ask whether abundance alone bounds its coclass and whether some element has centralizer of order at most \(p\) raised to a power depending only on abundance. Both bounds must be uniform over all primes. The source's small-abundance remarks retain their stated qualifications, including a missing boundary exception in its coclass claim.
\subsection{Sources}
\begin{itemize}
\item Local UnsolvedMath registry: permanent ID \(701073\), exact source problem code \texttt{catalogue entry}, title \emph{Kourovka Notebook Problem 14.89}; supplied upstream revision \texttt{d371ee917ac9647fbb06cd6fc529deec6f3dda82}. The supplied status is \texttt{open}; this is a registry label, not an independently verified assessment.
\item \emph{The Kourovka Notebook (Unsolved Problems in Group Theory)}, original problem 14.89, 14th issue (1999), as represented in the supplied record by the 21st issue edition (2026), arXiv identifier \texttt{1401.0300v48}. The collection metadata names E. I. Khukhro and V. D. Mazurov as editors.
\item Attribution is inconsistent within the supplied record: the background and proposer metadata name G. Fern\'andez--Alcober, with proposed year 1999, while the mathematical statement begins with attribution to E. A. O'Brien and A. Shalev. Both attributions are retained pending source review.
\item P. Hall: the conjugacy-class formula is attributed to a theorem of Hall in the supplied statement, without a specific bibliographic citation.
\item A. Jaikin-Zapirain, \emph{J. Group Theory} \textbf{3}, no. 3 (2000), 225--231. The supplied statement attributes the order bound depending on \(p\) and abundance to this reference; no article title is supplied.
\end{itemize}
\end{document}
